% Nouveau EGA V (cote 25): archive cote 25, site page 15 (PDF page 16).
% typescript. Draft compared: none.
% AI-transcribed by Claude Opus 5.5 (claude-opus-5-5, adaptive thinking, effort high), Message Batches API, batch ega25b (msgbatch_01A7n5YuXubUPNK1zsepi8RG), 2026-11-13; not yet checked by eye.
\scanpage{25/15}
\mspage{12}

\uline{Corollaire} 2.17.   Soit Z l'ensemble des points de X en lesquels F
n'est pas séparable sur k (réf). Alors  $G_{\hand{\unsure{y|\eta}}}$ est séparable sur $k(\hand{\unsure{y|\eta}})$ sss
Z est fini, En particulier, si F est séparable sur k, $G_{\hand{\unsure{y|\eta}}}$ est séparable
sur $k(\hand{\unsure{y|\eta}})$.

\margin{\unsure{bb pour|à mettre} dans l'énoncé 2.12.}\uline{Remarques} 2.18.  Le résultat clef  2.12. \inserted{(et ses conséquences \unsure{2.13.|2.11.} à 2.17.)} devient faux si on abandonne
l'hypothèse que f est une immer\unsure{z|s}ion, comme  on voit par exemple en
prenant pour  X un revêtement purement inséparable de P.  \xout{Cxxx} Cepen-%
dant \xout{lexrésultatxsuivantxx(quixauraitxdûxvenirxplûtôtxenxconclusion}
\xout{àx2x11xx)xxxxxSixxX} \xout{\unsure{Propxx}} si  k est \uline{de caractéristique nulle} ,  2.12. à
2.1\hand{7}. sont valables sans supposer f une immersion.  En effet, il suffit
de le vérifier pour 2.12. , et  cela résulte alors de 2.11. et du fait
que \xout{xxxxx} pour un préschéma algébrique en c ar 0,  lisse = régulier.

\deleted{Les propriétés envisagées jusqu'à maintenant étaient de nature
essentiellement locale. Nous allons terminer ce n° en étudiant les
propriétés globales}

\over{\unsure{3}}{2}.3.   \uline{Section hyperplane générique}: \uline{irréductibilité et connexité géomét.}
\uline{Théorème} 3.1. (Bertini-Zariski).  Supposons dim $f(X) \hand{\geq} 2$; et que X
soit \xout{géométriquementxintègrex(resp.x} géométriquement irréductible\deleted{)}. Alors
la section hyperplane générique $Y_{\hand{\unsure{y|\eta}}}$ a la même propriété.

\deleted{\xout{x}(Compte tenu que X séparable \hspace{3em} Y séparable, les deux énoncés
de 3.1. sont en fait équivalents: \xout{Onxpeutxdxailleursxsupposerx}  le
deuxième implique le premier, et inversement comme on voit en se rame-%
nant au cas k alg. clos.  Compte tenu de 2.1. et exmplicitant
la construction du corps des fonctions de Y \hspace{2em} comme extension de k( \hspace{1em} )
en termes de l'extension K de k et des éléments $t_1,\ldots,t_n$ de K images
d'un système de coordonnées affines convenable de}

\mspage{15}
